\section{Examples for a new task}

Use this file as a starting point for the next section of your report. The
examples below follow the style of the existing tasks. Rename the file,
replace the sample data, and remove any examples you do not need.

\subsection{Text and structure}
Split ordinary text into paragraphs and organize the content into subsections.
For example, introduce the model, present the simulation results, and then
discuss them.

Use a standard list for bullet points:
\begin{itemize}
    \item first item;
    \item second item;
    \item third item.
\end{itemize}

A numbered list:
\begin{enumerate}
    \item define the model;
    \item perform the calculation;
    \item analyze the result.
\end{enumerate}

\subsection{Equations and systems of equations}
Use the \texttt{equation} environment to add a number and a label to an equation:
\begin{equation}
    \dot{x} = Ax + Bu, \qquad y = Cx,
    \label{eq:taskN_model}
\end{equation}
You can then refer to it in the text: the model is defined by
equation~\eqref{eq:taskN_model}.

Use the \texttt{cases} environment for a system of equations:
\begin{equation}
    \begin{cases}
        \dot{x} = Ax + Bu, \\
        y = Cx.
    \end{cases}
\end{equation}

\subsection{Matrices and vectors}
Use \texttt{bmatrix} for an ordinary matrix:
\begin{equation}
    A = \begin{bmatrix}
        0 & 1 \\
        -2 & -3
    \end{bmatrix}.
\end{equation}

Use the row vector shortcut and separate entries with \texttt{\&}:
\begin{equation}
    x_0 = \rowvec{1 & 2 & 3}.
\end{equation}

For a column vector, separate entries with \texttt{\textbackslash\textbackslash}:
\begin{equation}
    b = \colvec{1 \\ -2 \\ 3}.
\end{equation}

The commands \texttt{\textbackslash rowvec} and \texttt{\textbackslash colvec}
create a matrix in square brackets and accept any number of entries.

\subsection{Results table}
\begin{table}[ht!]
    \centering
    \begin{tabular}{|c|c|c|}
        \hline
        Parameter & Value & Unit \\
        \hline
        $T$ & 3 & s \\
        $\alpha$ & 0.5 & -- \\
        \hline
    \end{tabular}
    \caption{Example results table}
    \label{tab:taskN_example}
\end{table}
Refer to the table using its label: table~\ref{tab:taskN_example}.

\subsection{Adding an image}
This shortcut inserts an image at the full text width:
\begin{verbatim}
\simpleimagepaste{media/plots/example.png}
\end{verbatim}

To add a caption and a reference label, supply them in square brackets:
\begin{verbatim}
\simpleimagepaste[Signal comparison][fig:taskN_comparison]
    {media/plots/example.png}
\end{verbatim}
Refer to this figure in the text with
\verb|\ref{fig:taskN_comparison}|.
You can also supply a caption without a label:
\begin{verbatim}
\simpleimagepaste[Output signal plot]{media/plots/example.png}
\end{verbatim}

To make an image narrower than the text, specify its width explicitly:
\begin{verbatim}
\imageatwidth[System diagram][fig:taskN_scheme]
    {0.7\textwidth}{media/scheme1.png}
\end{verbatim}

Paths are relative to the project root directory. The caption and label are
optional; omitting the caption inserts only the image.

% Example call (uncomment when the image file exists):
% \simpleimagepaste[Example plot][fig:taskN_plot]{media/plots/example.png}

\subsection{Cross-references}
Use \texttt{\textbackslash label} and \texttt{\textbackslash ref} for references
within the document, or \texttt{\textbackslash eqref} for equations.
For example: equation~\eqref{eq:taskN_model} defines the sample model.
